\chapter{Conclusion and Future Directions}
\label{ch:conclusion}

This dissertation extends the classical framework of statistical learning to settings where the data-generating process is influenced by strategic agents or where which variables are jointly observed departs from the traditional fixed-feature baseline. Together, these settings show that the performance of a learning system depends not only on the underlying environment, but also on the strategic incentives present in that environment and the structural composition of the observations through which it is perceived. Modern data-driven decision-making systems must therefore account for both strategic incentives and observational structure as integral components of the learning problem.

Across the strategic learning problems considered in this dissertation, a recurring theme is that strategic behavior need not always be viewed as an obstacle to be eliminated. When explicitly modeled, strategic behavior can reveal useful private information and induce structural properties that simplify the underlying decision problem. For example, in strategic feature revelation (Chapter~\ref{ch:strat_rev}), the sender's voluntary disclosure decision conveys information that, under non-erasure conditions, strictly improves predictive accuracy relative to the non-strategic baseline. Similarly, in IPO regulation (Chapter~\ref{ch:ipo}), permitting a bounded degree of presentation flexibility can align investor participation and strictly improve traded social welfare relative to complete prohibition or unregulated manipulation.

In the learning from partitioned observations problem (Chapter~\ref{ch:side_info}), which represents a structural rather than strategic departure, predictive performance depends fundamentally on which variables are observed simultaneously. When statistical dependencies between features and targets are never jointly observed during training, fundamental limitations arise on achievable risk. Within these limitations, the choice of imputation strategy governs learning performance. We show that, under specified environmental conditions, the regularizer that minimizes the theoretical bound on excess risk follows a three-candidate rule, settling at an interior value or taking one of two endpoints depending on the signal-to-imputation ratio.

The models studied throughout this dissertation are intentionally kept analytically tractable to expose transparent economic and statistical mechanisms. Although these stylized assumptions necessarily scope the immediate domain of application, they enable clean structural characterizations that provide useful principles for more complex systems.

Several directions remain open for future research. Beyond extending the proposed models to richer valuation and cost structures, an important direction is to explore non-stationary environments where agent preferences or observation partitions evolve dynamically over time. Such investigations naturally lie at the intersection of statistical learning and algorithmic game theory, with the goal of developing algorithms that are statistically efficient, robust to manipulation, and aware of observational constraints.

In conclusion, this dissertation demonstrates that statistical learning can be rigorously extended beyond its classical assumptions by explicitly modeling strategic incentives and partitioned observation structures. Doing so not only clarifies the fundamental limitations imposed by these departures, but also reveals new opportunities for designing decision systems that effectively exploit the underlying game-theoretic and structural properties of the problem.
